\section{Математическая постановка}
\label{sec:problem}

\subsection{Геометрия и история среды}
\draft{Основа модели --- постановка организаторов~\cite{problem2026} и принятые
в бассейновом моделировании зависимости~\cite{hantschel2009}.
Область $\Omega(t)=\{(x,z):0<x<W,\ z_{\mathrm{surf}}(x,t)<z<L(x,t)\}$
задана конфигурациями на границах эпох. Ось $z$ направлена вниз;
$t$ растёт от начала моделирования, геологический возраст $\tau$ уменьшается до нуля.
На рис.~\ref{fig:geometry} показаны три конфигурации и расположение наблюдений.}

\reportfigure{01_geometry_history.png}
  {Геометрия обратного набора при возрасте 200, 100 и 0 млн лет.
  Цвет обозначает литотип; на современной конфигурации показаны пять скважин и измерения T и Ro.}
  {fig:geometry}

\draft{Внутри эпохи контакты интерполируются по вертикали между заданными
конфигурациями. Принято постепенное отложение нового слоя:
его мощность линейно растёт от нуля до $h_{\mathrm{end}}(x)$.
Для конечной относительной глубины $a\in[0,1]$ положение материала и момент
его отложения определяются формулами:}
\begin{gather}
 \theta=\frac{t-t_k}{t_{k+1}-t_k},\qquad h(x,t)=\theta h_{\mathrm{end}}(x),\\
 z(a,t)=z_{\mathrm{surf}}(x,t)+h_{\mathrm{end}}(x)\max(0,\theta+a-1),\\
 t_{\mathrm{dep}}=t_k+(1-a)(t_{k+1}-t_k).
 \label{eq:deposition}
\end{gather}
\draft{Скорость $\boldsymbol v$ определяется движением материала.
Разломные перестройки происходят в заданные моменты на границах эпох.}

\subsection{Теплофизические свойства}
\draft{Свойства среды определяются свойствами скелета, флюида и пористостью.
Температура в формулах задаётся в К; $T_C=T-273{,}15$ --- в градусах Цельсия.}
\begin{gather}
 \lambda=\lambda_s(T)^{1-\varphi}\lambda_f^{\varphi},\qquad
 \rho c=(1-\varphi)\rho_s c_s(T)+\varphi\rho_f c_f,\qquad
 A=(1-\varphi)A_s,\\
 c_s(T)=c_{20}\left(0{,}953+2{,}29\cdot10^{-3}T_C
              -2{,}835\cdot10^{-6}T_C^2+1{,}191\cdot10^{-9}T_C^3\right),\\
 \lambda_s(T)=1{,}84+358(1{,}0227\lambda_{20}-1{,}882)
                            \left(\frac1T-0{,}00068\right).
\end{gather}
\draft{Зависимости теплоёмкости и теплопроводности --- Waples и Sekiguchi
соответственно~\cite{waples2004,sekiguchi1984}. В слое пористость линейна по
относительной глубине $a=(z-z_{\mathrm{top}})/(z_{\mathrm{bot}}-z_{\mathrm{top}})$:
$\varphi=(1-a)\varphi_{\mathrm{top}}+a\varphi_{\mathrm{bot}}$.
Граничные значения пористости линейно интерполируются по времени.}

\subsection{Теплоперенос и условия}
\draft{Температура удовлетворяет уравнению теплопереноса в движущемся материале.
Скорость $\boldsymbol v$ определяется заданной геометрической историей.
На внутренних контактах непрерывны температура и нормальный тепловой поток.}
\begin{gather}
 (\rho c)\left(\frac{\partial T}{\partial t}+\boldsymbol v\cdot\nabla T\right)
   =\nabla\cdot(\lambda\nabla T)+A,
 \label{eq:heat}\\
 T\big|_{z=z_{\mathrm{surf}}}=T_{\mathrm{surf}},\qquad
 \lambda\frac{\partial T}{\partial n}\bigg|_{z=L}=q_{\mathrm{basal}},\qquad
 \frac{\partial T}{\partial x}\bigg|_{x=0,W}=0,\\
 [T]=0,\qquad[\lambda\nabla T\cdot\boldsymbol n]=0.
\end{gather}

\draft{Начальное поле $T_0$ --- стационарное решение с начальными свойствами и
теми же граничными условиями при $t=0$.}
\begin{equation}
 -\nabla\cdot\bigl(\lambda(x,z,0,T_0)\nabla T_0\bigr)=A(x,z,0),
 \qquad T(x,z,0)=T_0(x,z).
\end{equation}

\subsection{Зрелость витринита}
\draft{Модель EASY\%Ro~\cite{sweeney1990} связывает зрелость с температурной
историей материальной точки $\boldsymbol\xi(s)$ от момента её отложения
$t_{\mathrm{dep}}$. Кинетические константы взяты из данных организаторов.
Для материала начальной конфигурации принято $t_{\mathrm{dep}}=0$.}
\begin{gather}
 I_i(t)=\int_{t_{\mathrm{dep}}}^{t}
 A_i\exp\!\left(-\frac{E_i}{R_gT(\boldsymbol\xi(s),s)}\right)\dd s,
 \qquad I_i(t_{\mathrm{dep}})=0,\\
 F=\sum_{i=1}^{N_R}f_i(1-e^{-I_i}),\qquad
 \Ro=\exp(-1{,}6+3{,}7F)\quad[\%].
\end{gather}

\subsection{Обратная задача}
\draft{Определяется только пространственно постоянный нижний поток $q(t)$,
монотонно убывающий во времени, в пределах 50--100 мВт/м$^2$.
Погрешности измерений приняты равными 2\% от заданных T в градусах Цельсия
и Ro. Выбрана нормировка невязок по этим погрешностям; регуляризации нет.}
\begin{gather}
 r_j^T=\frac{T_{C,j}^{\mathrm{calc}}-T_{C,j}^{\mathrm{obs}}}
                   {0{,}02|T_{C,j}^{\mathrm{obs}}|},\qquad
 r_j^{\Ro}=\frac{\Ro_j^{\mathrm{calc}}-\Ro_j^{\mathrm{obs}}}
                   {0{,}02|\Ro_j^{\mathrm{obs}}|},\\
 \min J,\qquad J=J_T+J_{\Ro}
       =\sum_j(r_j^T)^2+\sum_j(r_j^{\Ro})^2.
 \label{eq:objective}
\end{gather}
\draft{Согласование с погрешностью означает $|r_j|\leq1$ для каждого наблюдения.
Веса --- выбранная настройка метода; качество T и Ro также оценивается отдельно.}
